\section{Greek}

The Greek cosmos is genealogical. It is not built by a craftsman. It is
generated through a chain of divine generations and settles under \textbf{Zeus}.
The base text is \textbf{Hesiod's Theogony}, filled out by \textbf{Works and
Days} and later retold for Rome by \textbf{Ovid}.

\subsection{The Creation Model}

In the beginning there was only a \textbf{yawning gap}, an opening with nothing
in it. No maker stood behind it. It simply came first. Out of the gap, on their
own and begotten by no one, the first powers came into being. The \textbf{broad
earth} appeared, the seat of all things. The \textbf{misty depths} opened far
below. And \textbf{desire} appeared too, the force of joining that lets
everything else come together and breed.

From the earth, alone and without a mate, came the \textbf{starry sky},
stretched over her as her equal and her cover. Earth and sky were the
\textbf{first parents}. They brought forth a generation of children. The sky
hated his offspring and pressed them back down into the earth so none could be
born. The earth, in pain, made a weapon and asked her children to act. The
youngest and boldest ambushed his father and cut him down, taking his place as
the new ruler.

The new ruler feared the same fate. A warning said one of his own children would
unseat him, so he \textbf{swallowed each child} as it was born. His wife hid the
last one away and gave him a stone to swallow instead. The hidden child grew in
secret, returned, and forced his father to bring up the swallowed children. A
long \textbf{war between the generations} followed. The young gods won. They
cast the old ones down into the depths.

The youngest ruler then ordered the world. He and his two brothers
\textbf{divided the realms}. One took the sky, one took the sea, one took the
world of the dead below. The earth they held in common. The cosmos settled under
the rule of the youngest.

\subsection{The Model with Its Names}

\subsubsection{Chaos and the First Powers}

First of all came \textbf{Chaos}. The word means a \textbf{gap}, a yawning
opening, a void. It is no god and it acts on nothing. It is only the first state
from which the rest appears. After it the first powers arose, each on its own.

\begin{longtable}{L{0.22\linewidth} L{0.18\linewidth} L{0.42\linewidth}}
\caption{The Greek primordials, the first state and the powers born from it.}\\
\toprule
\textbf{Power} & \textbf{Greek} & \textbf{What it is} \\
\midrule
Chaos & Khaos & The yawning gap, first of all \\
Gaia & Gaia & The broad earth, firm seat of all \\
Tartarus & Tartaros & The misty pit deep below the earth \\
Eros & Eros & Desire, the joiner that lets the powers breed \\
\bottomrule
\end{longtable}

From Chaos came darkness (\textbf{Erebus}) and night (\textbf{Nyx}), and from
their union the bright upper air and the day. From \textbf{Gaia}, alone, came
the starry sky (\textbf{Ouranos}), the mountains, and the barren sea. \textbf{Nyx}
bore a brood of dark powers by herself, among them \textbf{Death}, \textbf{Sleep},
and the \textbf{Fates}.

\subsubsection{Ouranos, the Titans, and the First Succession}

\textbf{Gaia} lay with her own son \textbf{Ouranos}. They brought forth the
\textbf{Titans}, the elder gods, along with the one-eyed \textbf{Cyclopes} and the
hundred-handed \textbf{Hekatoncheires}. The youngest and most powerful Titan was
\textbf{Kronos}.

Ouranos hated his children and pushed them back into the earth so none could be
born. Gaia made a great \textbf{sickle of adamant} and asked her sons to strike.
Only Kronos was willing. He lay in ambush, and when Ouranos came down over Gaia at
night, Kronos \textbf{cut off his father} and threw the parts away. From the blood
that fell on the earth came the \textbf{Furies}, the \textbf{Giants}, and the
ash-tree nymphs. From the parts cast into the sea rose foam, and from the foam came
\textbf{Aphrodite}. This was the \textbf{first succession}. Sky was unseated by his
son.

\subsubsection{Kronos, Zeus, and the Titanomachy}

\textbf{Kronos} now ruled and married his sister \textbf{Rhea}. A warning said one
of his own children would overthrow him, so he \textbf{swallowed each child} as it
was born.

\begin{longtable}{L{0.22\linewidth} L{0.18\linewidth} L{0.42\linewidth}}
\caption{The children of Kronos and Rhea, swallowed in turn.}\\
\toprule
\textbf{Child} & \textbf{Greek} & \textbf{Later domain} \\
\midrule
Hestia & Hestia & Hearth \\
Demeter & Demeter & Grain \\
Hera & Hera & Marriage \\
Hades & Haides & The dead \\
Poseidon & Poseidon & The sea \\
Zeus & Zeus & The sky, the king \\
\bottomrule
\end{longtable}

Rhea hid the last child, \textbf{Zeus}, in a cave on \textbf{Crete} and gave
Kronos a \textbf{stone wrapped in cloth} to swallow in his place. Grown in secret,
Zeus forced Kronos to vomit up the swallowed children and the stone. This was the
\textbf{second succession}.

Then came the ten-year war of the gods against the
Titans (the \textbf{Titanomachy}). The Olympians fought from Mount \textbf{Olympus}, the Titans from Mount
\textbf{Othrys}. Zeus freed the Hekatoncheires and the Cyclopes from the earth.
The Cyclopes gave him the \textbf{thunderbolt}, and the hundred-handed buried the
Titans under a storm of rocks. The defeated Titans were cast into
\textbf{Tartarus} and bound. After Zeus also crushed the monster
\textbf{Typhoeus}, his rule was secure.

\subsubsection{The Division of the Cosmos by Lot}

The three brothers divided the world by \textbf{lot}, a casting of lots rather
than by birthright.

\begin{longtable}{L{0.22\linewidth} L{0.18\linewidth} L{0.42\linewidth}}
\caption{The division of the cosmos among the three brothers.}\\
\toprule
\textbf{God} & \textbf{Greek} & \textbf{Realm won by lot} \\
\midrule
Zeus & Zeus & The sky and the bright air \\
Poseidon & Poseidon & The sea \\
Hades & Haides & The underworld, the gloom \\
\bottomrule
\end{longtable}

The earth and Olympus were held in common by all three. Zeus stood first by the
lot of the sky.

\subsubsection{The Ages of Man}

From \textbf{Works and Days}, Hesiod describes five races of mortals, a steady
decline broken once by the heroes.

\begin{longtable}{L{0.16\linewidth} L{0.44\linewidth} L{0.24\linewidth}}
\caption{The five ages of man, a decline broken by the heroic age.}\\
\toprule
\textbf{Age} & \textbf{Character} & \textbf{End} \\
\midrule
Golden & Under Kronos, no toil, no grief, ease and feasting & Died as if asleep, became good spirits \\
Silver & Long childish, foolish, would not honor the gods & Zeus hid them away \\
Bronze & Terrible and strong, bronze arms, loved war & Destroyed by their own hands \\
Heroic & The demigods and heroes of Thebes and Troy, juster than bronze & Some dwell in the Isles of the Blessed \\
Iron & The present age, toil and sorrow, hard cares & Will worsen until Zeus ends it \\
\bottomrule
\end{longtable}

\subsubsection{Prometheus and Pandora}

The Titan \textbf{Prometheus} sided with men against Zeus. At a sacrifice he
tricked Zeus into choosing the bare bones hidden in shining fat, leaving the good
meat for men. In response Zeus \textbf{hid fire} from mortals. Prometheus
\textbf{stole fire} in a fennel stalk and gave it back. Zeus bound him to a rock
where an eagle ate his liver each day, the liver growing back each night.

To punish men for the fire, Zeus had a woman made. This was \textbf{Pandora}, the
first woman. \textbf{Hephaestus} shaped her from earth and water, \textbf{Athena}
dressed her, \textbf{Aphrodite} gave her grace, and \textbf{Hermes} gave her a
deceitful nature and a voice. She carried a \textbf{jar}, often called a box. She
opened it and released every \textbf{evil and disease} into the world. Only
\textbf{Hope} stayed inside under the rim. Woman was sent as the price of fire.

\subsubsection{The Olympians}

The twelve great gods of Zeus's generation dwell on Mount \textbf{Olympus}.

\begin{longtable}{L{0.22\linewidth} L{0.50\linewidth}}
\caption{The twelve Olympians. Hades is a brother of Zeus but dwells below and is not counted among them.}\\
\toprule
\textbf{God} & \textbf{Domain} \\
\midrule
Zeus & Sky, king of gods, oaths, hospitality \\
Hera & Marriage, queens, women \\
Poseidon & Sea, earthquakes, horses \\
Demeter & Grain, harvest, the seasons \\
Athena & Wisdom, craft, just war \\
Apollo & Sun, prophecy, music, healing \\
Artemis & Hunt, wild beasts, the moon \\
Ares & War, slaughter, courage \\
Aphrodite & Love, beauty, desire \\
Hephaestus & Fire, the forge, smithing \\
Hermes & Messages, travel, trade, guide of souls \\
Hestia or Dionysus & Hearth and home, or wine and ecstasy \\
\bottomrule
\end{longtable}

\subsubsection{The Underworld and the Soul}

The realm of \textbf{Hades} and his queen \textbf{Persephone} lies below. Souls
cross into it and are sorted. Five rivers wind through it. The \textbf{Styx} of
hatred, the \textbf{Acheron} of woe, the \textbf{Lethe} of forgetfulness, the
\textbf{Phlegethon} of fire, and the \textbf{Cocytus} of wailing.

The dead reach it by crossing the river. The ferryman \textbf{Charon} rows them
over, and they must pay him a \textbf{coin} placed in the mouth at burial. The
unburied wander the near shore. \textbf{Cerberus}, the three-headed dog, guards
the gate so none may leave. Three dead kings sit as \textbf{judges}.
\textbf{Minos} casts the final vote, \textbf{Rhadamanthus} judges the souls of the
East, and \textbf{Aeacus} judges those of the West.

\begin{longtable}{L{0.30\linewidth} L{0.46\linewidth}}
\caption{The three regions of the dead.}\\
\toprule
\textbf{Region} & \textbf{Who goes there} \\
\midrule
Elysium, the Isles of the Blessed & The blessed, the heroes, the initiated \\
The Asphodel Meadows & The ordinary dead, neither good nor wicked \\
Tartarus & The wicked and the punished Titans, the deep pit \\
\bottomrule
\end{longtable}

The named punished in Tartarus include \textbf{Sisyphus} with his rolling stone,
\textbf{Tantalus} with his receding food and water, and \textbf{Ixion} on his
burning wheel.

\subsubsection{The Fates}

Above the gods stand the powers of \textbf{fate}. Even Zeus is subject to them.
The three \textbf{Moirai} spin, measure, and cut the thread of each life.
\textbf{Clotho} spins it, \textbf{Lachesis} measures its length, and
\textbf{Atropos} cuts it. Behind them stands \textbf{Ananke}, necessity itself, a
binding the gods cannot break.

\subsubsection{The Mystery Cults}

Beyond the public worship lay the \textbf{mysteries}, secret initiations that
promised the initiate a \textbf{better lot after death}. They are the esoteric
part of Greek religion.

The \textbf{Eleusinian Mysteries} of \textbf{Demeter} and \textbf{Persephone}
were the most prominent. \textbf{Hades} seizes \textbf{Persephone} and takes her
below. \textbf{Demeter} wanders in grief and lets no crops grow, until Zeus sends
\textbf{Hermes} to bring the daughter back. But she has eaten \textbf{pomegranate
seeds} below, so she must return for part of each year. Her yearly descent and
return account for the \textbf{seasons}. The grain that dies in the ground and
rises again is the sign, and the initiate who has \textbf{seen} the rites is
promised a better afterlife.

The \textbf{Orphic} and \textbf{Dionysian} mysteries went further. Their central
myth tells how the child \textbf{Dionysus Zagreus} is torn apart and eaten by the
Titans. Zeus blasts the Titans to ash, and from that ash humanity is born. So each
person holds \textbf{two natures}, the earthly \textbf{Titanic} part and the
divine \textbf{Dionysian} spark. The Orphic life aims to purify away the Titanic
and free the divine soul.

This rests on the teaching of the
\textbf{transmigration of souls} (\textbf{metempsychosis}). The soul is immortal and divine, fallen into the
body, which is its \textbf{tomb}. At death the unpurified soul drinks
\textbf{Lethe}, forgets, and is reborn. Through purification and the mysteries it
can break the cycle of rebirth and return to its divine origin. Small
\textbf{gold tablets} buried with initiates carried the password for the soul. It
must avoid the spring of Forgetfulness, drink instead from the spring of
\textbf{Memory}, and declare, \textbf{I am a child of Earth and starry Sky, but my
race is of Heaven}.

\subsection{Comparison}

The Greek myth states the pattern of this collection in concrete terms. The
\textbf{one goes out into the many and returns}. Out of the single \textbf{gap}
the powers come forth on their own, then breed generation after generation into
the full set of gods and the ordered world. That is the going-out, told as
begetting and war rather than emanation. The return is not in the public myth but
in the \textbf{mysteries}. There the fallen soul, scattered into body after body,
drinks from Memory, breaks the cycle of rebirth, and returns to its divine source.
The base models in this collection name the source and map every stage of its
descent and repair. The Greek myth leaves the source as a faceless gap and tells
the rest as a family narrative. The same arc runs underneath. The many come out
of a first state, and the soul finds its way back to where it began.
